\documentclass[14pt,a4paper]{extarticle}

\usepackage[utf8]{inputenc}
\usepackage[T2A]{fontenc}
\usepackage[russian,english]{babel}
\usepackage{geometry}
\usepackage{setspace}
\usepackage{titlesec}
\usepackage{tocloft}
\usepackage{graphicx}
\usepackage{float}
\usepackage{booktabs}
\usepackage{longtable}
\usepackage{array}
\usepackage{hyperref}
\usepackage{tikz}
\usetikzlibrary{positioning,shapes.geometric,arrows.meta,calc}
\usepackage{listings}
\usepackage{xcolor}

\geometry{
  left=30mm,
  right=15mm,
  top=20mm,
  bottom=20mm
}

\onehalfspacing
\setlength{\parindent}{1.25cm}
\setlength{\parskip}{0pt}

\titleformat{\section}{\normalfont\bfseries\large}{\thesection}{1em}{}
\titleformat{\subsection}{\normalfont\bfseries}{\thesubsection}{1em}{}

\renewcommand{\contentsname}{СОДЕРЖАНИЕ}
\makeatletter
\renewcommand\@biblabel[1]{#1.}
\makeatother

\lstset{
  basicstyle=\ttfamily\small,
  breaklines=true,
  frame=single,
  backgroundcolor=\color{gray!5},
  columns=fullflexible
}

% --- Metadata ---
\newcommand{\University}{Министерство науки и высшего образования Российской Федерации\\
Федеральное государственное бюджетное образовательное учреждение высшего образования\\
«Московский авиационный институт (национальный исследовательский университет)»}
\newcommand{\Department}{Кафедра №806}
\newcommand{\Supervisor}{Крылов С.С.}
\newcommand{\Advisor}{Кубраков Д.В.}
\newcommand{\Student}{ФИО студента}
\newcommand{\Group}{М8О-000Б-00}
\newcommand{\City}{Москва}
\newcommand{\Year}{2026}
\newcommand{\ReportTitle}{Разработка индексера блокчейна Bitcoin с REST API и модульной архитектурой}

\begin{document}

% --- Title page ---
\pagenumbering{gobble}
\begin{center}
\University
\end{center}

\vspace{1.5cm}

\begin{flushright}
\textbf{УТВЕРЖДАЮ}\\
Заведующий \Department\\
\Supervisor\\
«\hspace{1.5cm}»\hspace{2cm}\Year\,г.
\end{flushright}

\vspace{1.2cm}

\begin{center}
\textbf{ОТЧЁТ}\\
О НАУЧНО-ИССЛЕДОВАТЕЛЬСКОЙ РАБОТЕ\\
\vspace{0.4cm}
ПОЯСНИТЕЛЬНАЯ ЗАПИСКА\\
по комплексной лабораторной работе\\
по дисциплине «Программная инженерия»\\
по теме:\\
\ReportTitle
\end{center}

\vspace{1.0cm}

\begin{tabular}{p{0.6\textwidth} p{0.35\textwidth}}
Ассистент кафедры & \Advisor \\
Студент группы \Group & \Student \\
\end{tabular}

\vfill

\begin{center}
\City\, \Year
\end{center}

\newpage

% --- List of executors ---
\pagenumbering{arabic}
\setcounter{page}{2}
\begin{center}
\textbf{СПИСОК ИСПОЛНИТЕЛЕЙ}
\end{center}

\vspace{0.8cm}

\noindent Студент группы \Group \hfill \Student

\newpage

% --- Abstract ---
\begin{center}
\textbf{РЕФЕРАТ}
\end{center}

\noindent Пояснительная записка содержит 20 стр., 5 рис., 3 табл., 6 источн., 1 прил.

\noindent Ключевые слова: БЛОКЧЕЙН BITCOIN, ИНДЕКСАЦИЯ ДАННЫХ, ИНДЕКСЕР, BACKEND-СЕРВИС, RUST, POSTGRESQL, REST API, JSON-RPC, BITCOIN CORE, JOBS, MEMPOOL, REORG, UTXO, MTLS, BASIC AUTH, DOCKER COMPOSE, PROMETHEUS, STATELESS-АРХИТЕКТУРА, МОДУЛЬНЫЙ МОНОЛИТ.

\noindent Работа направлена на разработку и обоснование технических решений для построения backend-системы индексации блокчейна Bitcoin, обеспечивающей консистентное хранение on-chain данных и их предоставление внешним потребителям через защищенный программный интерфейс.

\noindent Объектом исследования является программный продукт «Bitcoin Blockchain Indexer» — backend-сервис на Rust \cite{src_rust}, получающий данные из Bitcoin Core по RPC, нормализующий их, сохраняющий в PostgreSQL \cite{src_postgres} и предоставляющий REST API для управления jobs и выдачи данных.

\noindent Предметом исследования являются методы и средства построения отказоустойчивого статeless-индексера: алгоритмы синхронизации блоков и транзакций, обработка reorg и mempool, организация модели данных и обеспечение защищенного взаимодействия компонентов.

\noindent Целью работы является проектирование и реализация MVP индексера блокчейна Bitcoin, обеспечивающего полную индексацию блоков и транзакций, поддержку reorg заданной глубины, мониторинг mempool, наблюдаемость и безопасное взаимодействие с узлом через mTLS и Basic Auth.

\noindent Для достижения цели решены следующие задачи: выполнен анализ требований и ограничений предметной области; спроектирована архитектура модульного монолита и структура конфигурации; разработана схема данных PostgreSQL и механизмы миграций; реализованы процессы индексации блоков, транзакций и mempool; реализовано API управления jobs и выдачи данных; обеспечены аутентификация и шифрование каналов; подготовлены метрики, журналирование и средства контейнерного развёртывания; разработана программа и методика испытаний.

\noindent В работе использованы методы системного и структурного анализа, моделирование архитектуры и потоков данных, проектирование реляционных структур хранения, методы модульного и интеграционного тестирования, а также экспериментальная проверка работоспособности на стенде docker-compose.

\noindent Результатом работы является реализованный MVP сервиса «Bitcoin Blockchain Indexer» с документированной архитектурой, реализованными ключевыми функциями индексации и подтвержденной работоспособностью по результатам испытаний.

\noindent Практическая значимость работы заключается в возможности использования разработанного решения как базового компонента аналитических и сервисных систем, работающих с данными блокчейна Bitcoin, а также как расширяемой основы для дальнейшего развития в рамках ВКР.

\newpage

% --- Contents ---
\tableofcontents
\newpage

% --- Definitions ---
\begin{center}
\textbf{ОПРЕДЕЛЕНИЯ}
\end{center}

\noindent В настоящем отчете применяются следующие термины с соответствующими определениями.

\noindent \textbf{Индексер} — backend-сервис, получающий данные блокчейна, нормализующий их и сохраняющий в БД с предоставлением API.

\noindent \textbf{Узел Bitcoin (Bitcoin Core)} — программный узел сети Bitcoin, предоставляющий доступ к данным блокчейна и mempool через JSON-RPC интерфейс.

\noindent \textbf{Каноническая цепь} — актуальная на текущий момент ветка блоков, признаваемая сетью основной.

\noindent \textbf{Reorg} — ситуация смены канонической цепи блоков, приводящая к откату ранее принятых блоков.

\noindent \textbf{Orphaned-блок (неканонический блок)} — блок, исключенный из канонической цепи в результате reorg.

\noindent \textbf{Mempool} — набор неподтвержденных транзакций узла, ожидающих включения в блок.

\noindent \textbf{On-chain данные} — данные подтвержденных транзакций и блоков, включенных в каноническую цепь.

\noindent \textbf{UTXO} — непотраченный выход транзакции.

\noindent \textbf{Job} — управляемый процесс индексации с собственным идентификатором, статусом и прогрессом выполнения.

\noindent \textbf{Жизненный цикл job} — допустимая последовательность состояний задания индексации (создание, запуск, пауза, возобновление, завершение/ошибка) с валидацией переходов.

\noindent \textbf{Tip цепи} — наибольшая известная высота канонической цепи блоков на узле.

\noindent \textbf{Идемпотентность индексации} — свойство, при котором повторная обработка одних и тех же входных данных не приводит к неконсистентным дубликатам или изменению итогового корректного состояния БД.

\noindent \textbf{Stateless-сервис} — сервис, который не хранит критичное бизнес-состояние в локальной памяти процесса между перезапусками, а использует внешние хранилища и конфигурацию.

\noindent \textbf{Модульный монолит} — архитектурный стиль, при котором система разрабатывается как единое приложение, но разделяется на изолированные модули с явными границами ответственности.

\noindent \textbf{Reverse proxy} — промежуточный сервис, через который выполняется доступ к целевому узлу; в проекте используется для контролируемого и защищенного RPC-взаимодействия.

\noindent \textbf{Наблюдаемость} — совокупность средств мониторинга и диагностики (метрики, логи, трассировка состояний), позволяющая оценивать работоспособность и качество функционирования сервиса.

\noindent \textbf{Миграции БД} — управляемые версии изменений схемы базы данных, применяемые для приведения структуры хранения к требуемому состоянию.

\newpage

% --- Abbreviations ---
\begin{center}
\textbf{ОБОЗНАЧЕНИЯ И СОКРАЩЕНИЯ}
\end{center}

\noindent ПО — программное обеспечение.

\noindent БД — база данных.

\noindent MVP — Minimum Viable Product, минимально жизнеспособный продукт.

\noindent RPC — Remote Procedure Call, JSON-RPC интерфейс Bitcoin Core.

\noindent JSON — JavaScript Object Notation, текстовый формат обмена данными.

\noindent JSON-RPC — протокол удаленного вызова процедур поверх JSON.

\noindent API — Application Programming Interface.

\noindent REST — Representational State Transfer.

\noindent HTTP — HyperText Transfer Protocol.

\noindent HTTPS — HTTP over TLS, защищенный протокол передачи данных.

\noindent SQL — Structured Query Language.

\noindent CLI — Command Line Interface.

\noindent YAML — YAML Ain't Markup Language, формат конфигурационных файлов.

\noindent TLS — Transport Layer Security.

\noindent mTLS — взаимная аутентификация по TLS сертификатам.

\noindent CI/CD — Continuous Integration / Continuous Delivery, практики непрерывной интеграции и поставки.

\noindent SLA — Service Level Agreement, соглашение об уровне сервиса.

\noindent URL — Uniform Resource Locator, адрес сетевого ресурса.

\noindent PK — Primary Key, первичный ключ таблицы.

\noindent FK — Foreign Key, внешний ключ таблицы.

\noindent UQ — Unique Constraint, ограничение уникальности.

\noindent ER — Entity-Relationship, модель «сущность--связь».

\noindent DFD — Data Flow Diagram, диаграмма потоков данных.

\noindent BPMN — Business Process Model and Notation.

\noindent UML — Unified Modeling Language.

\newpage

% --- Introduction ---
\section*{Введение}
\addcontentsline{toc}{section}{Введение}

\noindent Поставленная проблема состоит в том, что прямой доступ к данным блокчейна Bitcoin через узел не обеспечивает требуемую для прикладных систем скорость аналитических выборок, устойчивую воспроизводимость результатов и удобный интерфейс управления процессом обработки данных \cite{src_bitcoin_core}. В практической эксплуатации это проявляется как рост времени ответа на запросы, высокая нагрузка на RPC-интерфейс узла, сложность повторяемого формирования витрин данных и риск рассогласования данных при reorg и изменениях mempool.

\noindent На практике указанная проблема возникает в системах мониторинга транзакций, антифрод-аналитике, внутренней отчетности, сервисах уведомлений и продуктовых API, где требуется многократное чтение исторических и оперативных on-chain данных. Для таких сценариев необходим специализированный индексирующий слой, отделяющий контур чтения клиентских сервисов от контура получения первичных данных из узла.

\noindent Необходимость разработки собственного индексера обусловлена требованиями безопасности, управляемости и предсказуемости качества сервиса: данные должны храниться в контролируемой инфраструктуре, доступ к ним должен предоставляться через стабильный REST API-контракт, а критичные процессы индексации должны быть наблюдаемыми и проверяемыми. Аналогичный подход применяется на рынке: собственные платформы индексации и обработки блокчейн-данных развивают, в частности, Coinbase, Chainalysis, Blockstream и Binance.

\noindent Цель работы формулируется как разработка и описание MVP индексера блокчейна Bitcoin, который обеспечивает детерминированную индексацию блоков и транзакций, обработку reorg, мониторинг mempool, безопасную RPC-интеграцию и предоставление данных через защищенный API.

\noindent Исходными данными для разработки служат: техническое задание проекта; спецификация и документация RPC Bitcoin Core; архитектурные ограничения (модульный монолит, stateless); требования к безопасности (TLS/mTLS, Basic Auth); требования к хранению данных в PostgreSQL и контейнерному развёртыванию.

\noindent Разрабатываемое решение должно включать backend-сервис на Rust, конфигурируемый механизм jobs, подсистему индексации блоков/транзакций/адресных данных, обработку канонической и неканонической ветвей цепи, учет mempool-транзакций, REST API управления и выдачи данных, а также средства наблюдаемости (метрики и логи) и инфраструктурные артефакты docker-compose.

\noindent Метрологическое обеспечение работы рассматривается как система измерений и контроля качества, включающая формализацию показателей (пропускная способность, латентность индексации, доступность API, корректность данных), методику их вычисления, тестовые сценарии и правила приемки результатов. Проверка выполняется на фиксированных профилях нагрузки с сопоставлением состояния API, БД, логов и метрик Prometheus.

\noindent Актуальность темы обусловлена одновременным ростом объема блокчейн-данных и числа прикладных сервисов, использующих эти данные в операционных и аналитических процессах. В этих условиях создание собственного индексера является необходимым шагом для обеспечения независимости от внешних провайдеров, повышения скорости доступа к данным и обеспечения требуемого уровня надежности и безопасности.

\newpage

\section{Вербальное описание области, направления работы, места и роли продукта}

\subsection{Место, роль и назначение результата}
\noindent Результатом главы является вербальное описание предметной области и направления разработки индексера блокчейна Bitcoin. Место результата в структуре ПЗ -- формирование исходной содержательной рамки для последующих артефактов: инфологической модели, архитектурных решений, метрик качества и программы испытаний. Роль результата состоит в фиксации контекста применения продукта, границ системы, ожидаемых свойств и ограничений, которые должны сохраняться на всех этапах проектирования и реализации.

\subsubsection{Какие функции выполняет данное описание}
\noindent Вербальное описание выполняет следующие функции:
\begin{enumerate}
  \item \textbf{Контекстная функция.} Определяет проблемное поле и практические сценарии использования индексера.
  \item \textbf{Коммуникационная функция.} Формирует единый понятийный аппарат для разработчиков, заказчика и проверяющей стороны.
  \item \textbf{Трассировочная функция.} Обеспечивает связь между требованиями ТЗ и проектными решениями в последующих главах.
  \item \textbf{Ограничительная функция.} Фиксирует границы MVP и исключает расширение реализации за пределы согласованных требований.
  \item \textbf{Критериальная функция.} Задает основу для оценки качества полученного решения и проверки достижения цели работы.
\end{enumerate}

\subsection{Описание предметной области}
\noindent Предметная область охватывает процессы получения, нормализации, хранения и предоставления данных блокчейна Bitcoin для внутренних и внешних потребителей. Источником первичных данных выступает узел Bitcoin Core, доступный через RPC. Целевыми потребителями являются сервисы аналитики, мониторинга, отчетности и интеграционные компоненты, которым необходимы быстрые и детерминированные выборки по блокам, транзакциям, адресам и состояниям jobs.

\subsubsection{Современные подходы к индексации}
\noindent В современной практике применяются три базовых подхода: прямой доступ к узлу по RPC, использование сторонних публичных API и построение собственного индексирующего backend-слоя. Первый подход прост во внедрении, но плохо масштабируется при росте числа сложных выборок. Второй снижает затраты на сопровождение инфраструктуры, однако ограничивает контроль над безопасностью, SLA и полнотой данных. Третий подход требует больших начальных затрат, но обеспечивает контроль над моделью данных, предсказуемую производительность и возможность адаптации под внутренние бизнес-сценарии.

\subsubsection{Проблемы существующих решений}
\noindent Для прямого RPC-доступа характерны повышенная латентность комплексных запросов, высокая нагрузка на узел и ограниченная пригодность для повторяемой аналитики. При использовании внешних индексаторов возникают риски зависимости от третьей стороны, ограничения на глубину/полноту данных, вариативность контрактов API и дополнительные риски информационной безопасности. Также типичной проблемой является неполная обработка reorg/mempool-сценариев, что приводит к рассогласованию статусов данных в прикладных системах.

\subsubsection{Место и роль интеллектуальной системы}
\noindent В рамках данной работы роль интеллектуальной системы выполняет программный комплекс «Bitcoin Blockchain Indexer», который является связующим звеном между инфраструктурой блокчейн-узла и потребителями данных. Система предназначена для непрерывного управляемого цикла индексации и предоставления данных через единый API-контракт. Основные функции системы: получение данных из Bitcoin Core, нормализация и запись в PostgreSQL, управление жизненным циклом jobs, обработка reorg, учет mempool-транзакций, выдача данных и административных состояний через REST API, а также формирование метрик и логов для наблюдаемости. Целевые показатели включают полноту индексации, детерминизм результатов, допустимую латентность выборок, корректность обработки reorg и доступность API. Система взаимодействует с Bitcoin Core (через RPC proxy), PostgreSQL, клиентскими сервисами и контуром мониторинга Prometheus.

\subsection{Характеристика результата}
\noindent Полученный результат представляет собой целостный вербальный артефакт, в котором зафиксированы проблематика, границы решения, практическая роль системы, ключевые функции и критерии назначения продукта. Артефакт является достаточным для перехода к формализации сущностей и связей в инфологической модели, а также для обоснования архитектуры модульного монолита со stateless-поведением. Результат характеризуется прикладной направленностью: все формулировки ориентированы на проверяемость в проектировании, реализации и испытаниях MVP.

\subsection{Оценка качества артефакта}
\noindent Качество вербального артефакта оценивается по четырем критериям: полнота, непротиворечивость, трассируемость и проверяемость. Полнота подтверждается покрытием проблемы, контекста применения, функций и ограничений системы. Непротиворечивость обеспечивается согласованием формулировок с ТЗ и терминологическим аппаратом записки. Трассируемость выражена в прямой связи положений главы с последующими разделами по модели данных, архитектуре, метрикам и испытаниям. Проверяемость обеспечивается наличием целевых показателей, которые могут быть измерены в рамках программы и методики испытаний.

\newpage

\section{Инфологическая схема предметной области продукта}

\subsection{Место, роль и назначение результата}
\noindent В данной главе используется инфологическая модель ER в нотации Чена, которая определяет состав сущностей предметной области, их атрибуты и семантические связи до перехода к физической схеме хранения. Модель нужна для согласования требований главы 1 с архитектурой и БД-слоем: она фиксирует структуру данных индексера, устраняет неоднозначности трактовки терминов и задает основу для проектирования API-контрактов и правил целостности.

\subsection{Изложение артефакта}
\noindent Инфологическая схема представлена ER-диаграммой в нотации Чена, вынесенной в приложение Б (рисунок~\ref{fig:chen-er}). Схема объединяет функциональные данные блокчейн-индексации (блоки, транзакции, входы/выходы, UTXO, балансы) и сервисные данные управления процессом (jobs, адресные выборки jobs, состояние узла), обеспечивая единое представление предметной области для последующего проектирования.

\subsubsection{Описание сущностей}
\begin{longtable}{p{0.22\textwidth} p{0.72\textwidth}}
\caption{Сущности инфологической модели и их назначение} \\
\toprule
Сущность & Подробное описание \\
\midrule
\endfirsthead
\toprule
Сущность & Подробное описание \\
\midrule
\endhead
\texttt{blocks} & Сущность канонической и неканонической цепи. Хранит блок на определенной высоте, его хэш и ссылку на предыдущий блок, временную метку и статус каноничности. Используется как опорная сущность для трассировки состояния цепи и обработки reorg. \\
\texttt{transactions} & Сущность транзакций блокчейна и mempool. Для каждой транзакции фиксируются идентификатор, привязка к блоку (при подтверждении), статус и декодированное содержимое. Сущность связывает блоковый уровень и детальный уровень vin/vout-структур. \\
\texttt{tx\_outputs} & Сущность выходов транзакций. Описывает каждый выход через пару \texttt{txid/vout}, сумму в сатоши, тип скрипта, адрес и скриптовые данные. Выходы являются источником формирования UTXO и адресных балансов. \\
\texttt{tx\_inputs} & Сущность входов транзакций. Для каждого входа хранит ссылку на ранее созданный выход (\texttt{prev\_txid/prev\_vout}) и служебные параметры. Используется для фиксации расходования UTXO и построения цепочки происхождения средств. \\
\texttt{utxos\_current} & Сущность текущего UTXO-состояния. Представляет актуальный набор непотраченных выходов и их статус, включая факт последующего расходования. Применяется для быстрых адресных выборок и расчета балансов. \\
\texttt{address\_balance\_current} & Сущность текущего баланса адреса. Хранит последнее агрегированное состояние баланса для низколатентного чтения через API. Обновляется при изменении UTXO-состояния. \\
\texttt{address\_balance\_history} & Сущность исторических срезов баланса адреса по высотам блоков. Нужна для ретроспективной аналитики и воспроизводимых отчетов. \\
\texttt{jobs} & Сущность управляемых процессов индексации. Хранит идентификатор задания, режим работы, статус, прогресс, диагностические поля и снимок конфигурации. Используется для оркестрации жизненного цикла индексации. \\
\texttt{job\_addresses} & Сущность адресного фильтра для jobs режима \texttt{address\_list}. Хранит множество адресов, связанных с конкретным заданием. \\
\texttt{node\_health} & Сущность наблюдаемости RPC-узла. Содержит актуальные параметры доступности и состояния узла (время последнего контакта, tip, latency, статус), используемые для мониторинга и диагностики. \\
\bottomrule
\end{longtable}

\subsubsection{Описание связей}
\begin{longtable}{p{0.24\textwidth} p{0.14\textwidth} p{0.56\textwidth}}
\caption{Связи инфологической модели в нотации Чена} \\
\toprule
Связь & Кардинальность & Подробное описание \\
\midrule
\endfirsthead
\toprule
Связь & Кардинальность & Подробное описание \\
\midrule
\endhead
\texttt{contains} (\texttt{blocks}--\texttt{transactions}) & 1:N & Один блок содержит множество транзакций; каждая подтвержденная транзакция принадлежит одному блоку. Связь отражает базовую иерархию блокчейна. \\
\texttt{has\_output} (\texttt{transactions}--\texttt{tx\_outputs}) & 1:N & Одна транзакция формирует один или более выходов. Связь используется для последующего построения UTXO и адресной агрегации. \\
\texttt{has\_input} (\texttt{transactions}--\texttt{tx\_inputs}) & 1:N & Одна транзакция содержит один или более входов (кроме coinbase-специфики). Связь фиксирует структуру потребления ранее созданных выходов. \\
\texttt{spends} (\texttt{tx\_inputs}--\texttt{tx\_outputs}) & N:1 & Вход транзакции ссылается на конкретный ранее созданный выход. Для выхода допустимо не более одного факта расходования в канонической истории, что обеспечивает корректность UTXO-модели. \\
\texttt{forms\_utxo} (\texttt{tx\_outputs}--\texttt{utxos\_current}) & 1:0..1 & Не каждый выход попадает в текущее UTXO-состояние (например, уже потраченный), но каждая UTXO-запись происходит из конкретного выхода. \\
\texttt{updates\_current\_balance} (\texttt{utxos\_current}--\texttt{address\_balance\_current}) & N:1 & Множество UTXO одного адреса агрегируется в одну текущую балансовую запись. Связь реализует оптимизацию чтения адресного состояния. \\
\texttt{snapshots} (\texttt{address\_balance\_current}--\texttt{address\_balance\_history}) & 1:N & Для одной текущей балансовой записи формируется множество исторических срезов по высотам блоков. \\
\texttt{has\_address\_filter} (\texttt{jobs}--\texttt{job\_addresses}) & 1:N & Одно задание режима адресного мониторинга может включать множество адресов наблюдения; каждая запись адреса принадлежит одному job. \\
\texttt{monitors\_node} (\texttt{jobs}--\texttt{node\_health}) & N:1 & Несколько jobs могут использовать одно состояние узла; связь отражает зависимость рабочих процессов от доступности RPC-инфраструктуры. \\
\bottomrule
\end{longtable}

\subsection{Описание условных обозначений}
\noindent Для ER-диаграммы в нотации Чена используются следующие обозначения: прямоугольник -- сущность; ромб -- связь; овал -- атрибут; подчеркнутый атрибут -- ключевой атрибут сущности; подписи \texttt{1}, \texttt{N}, \texttt{0..1} у ребер -- кардинальность участия сущностей в связи. Именованные ромбы отражают семантику предметных отношений, а не физические внешние ключи. Используемая нотация основана на классическом подходе П. Чена \cite{src_chen}.

\subsection{Характеристика результата}
\noindent Сформирована инфологическая модель, содержащая 10 сущностей и 9 ключевых связей предметной области индексера. Такая структура позволяет одновременно покрыть контур индексации блокчейн-данных и контур операционного управления jobs. За счет разделения на транзакционный, агрегированный и сервисный слои достигаются: однозначная трассировка происхождения данных, возможность детерминированной обработки reorg/mempool-событий и поддержка низколатентных клиентских выборок.

\subsection{Оценка качества артефакта}
\noindent Качество ER-диаграммы определяется по следующим критериям.
\begin{enumerate}
  \item \textbf{Полнота предметного покрытия.} Оценивается наличием всех сущностей и связей, необходимых для реализации требований ТЗ и бизнес-правил главы 1.
  \item \textbf{Семантическая корректность.} Проверяется соответствие кардинальностей и названий связей реальной логике блокчейн-данных и жизненного цикла jobs.
  \item \textbf{Непротиворечивость.} Оценивается отсутствие конфликтов между определениями сущностей, их атрибутами и отношениями.
  \item \textbf{Трассируемость к проектным решениям.} Проверяется возможность прямого перехода от ER-модели к логической схеме БД, API-операциям и тестовым сценариям.
  \item \textbf{Расширяемость.} Оценивается возможность добавления новых сущностей и связей (например, дополнительных индексов и аналитических представлений) без разрушения текущей структуры.
  \item \textbf{Проверяемость.} Определяется возможностью формально проверить модель через ограничения целостности и сценарии испытаний.
\end{enumerate}

\newpage

\section{Модель функционирования организации или устройства внедрения}

\subsection{Место, роль и назначение результата}
\noindent В главе используется модель потоков данных DFD (Data Flow Diagram), позволяющая описать функционирование системы с точки зрения движения данных между внешними участниками, процессами и хранилищами. Модель определяет границы ответственности индексера в контуре эксплуатации, состав функциональных процессов и точки интеграции с узлом Bitcoin Core, базой данных и клиентскими сервисами. Назначение результата -- обеспечить согласование бизнес-правил из главы 1 и инфологической модели главы 2 с операционной логикой работы продукта.

\subsection{Изложение артефакта}
\noindent Набор DFD-диаграмм вынесен в приложение В: DFD уровня 0 (контекст), DFD уровня 1 (декомпозиция основных процессов) и DFD уровня 2 (детализация процессов индексации, администрирования и получения данных при отсутствии данных в индексере). Совокупно диаграммы обеспечивают прослеживаемость потоков данных от внешнего запроса до сохранения/выдачи результата, а также фиксируют поведение системы в штатных и граничных сценариях.

\subsubsection{DFD уровня 0 (контекстная диаграмма)}
\noindent Контекстная диаграмма (приложение В, DFD 0 уровня) показывает АС как единый процесс \texttt{P0 Bitcoin Blockchain Indexer} и отражает взаимодействие с внешними сущностями:
\begin{enumerate}
  \item \textbf{E1 Клиентские системы} -- источник управляющих и пользовательских запросов и получатель результатов обработки.
  Потоки данных:
  \begin{itemize}
    \item вход в АС: \texttt{API-запросы: данные, команды управления};
    \item выход из АС: \texttt{API-ответы: данные индекса}.
  \end{itemize}
  \item \textbf{E2 Bitcoin Core RPC Proxy} -- внешний источник первичных блокчейн-данных.
  Потоки данных:
  \begin{itemize}
    \item выход из АС: \texttt{JSON-RPC запросы к узлу};
    \item вход в АС: \texttt{Блоки, транзакции, mempool}.
  \end{itemize}
  \item \textbf{E3 Контур мониторинга} -- потребитель телеметрии и источник эксплуатационных сигналов.
  Потоки данных:
  \begin{itemize}
    \item выход из АС: \texttt{Метрики, логи};
    \item вход в АС: \texttt{Сигналы и алерты}.
  \end{itemize}
\end{enumerate}
\noindent Таким образом, на контекстном уровне фиксируются основные контуры: клиентский API-обмен, сбор первичных данных из узла и контур наблюдаемости.

\subsubsection{DFD уровня 1}
\noindent DFD 1 уровня (приложение В, DFD 1 уровня) раскрывает внутреннюю структуру индексера и включает три процесса:
\begin{enumerate}
  \item \textbf{P1 Индексация данных.} Формирует и поддерживает актуальный индекс блокчейн-данных.
  Потоки:
  \begin{itemize}
    \item вход: \texttt{Сигналы запуска/паузы/возобновления} (из \texttt{P2}), \texttt{Параметры активных jobs} (из \texttt{D2}), \texttt{Ответы RPC} (из \texttt{E2});
    \item выход: \texttt{RPC-запросы на блоки/транзакции/mempool} (в \texttt{E2}), \texttt{Нормализованные блоки/tx/utxo/балансы} (в \texttt{D1}), \texttt{Прогресс jobs, ошибки, служебные статусы} (в \texttt{D2}), \texttt{Метрики индексации и логи} (в \texttt{E3}).
  \end{itemize}
  \item \textbf{P2 Администрирование jobs.} Управляет жизненным циклом заданий и их параметрами.
  Потоки:
  \begin{itemize}
    \item вход: \texttt{Команды управления jobs} (из \texttt{E1}), \texttt{Текущие состояния и параметры jobs} (из \texttt{D2});
    \item выход: \texttt{Статусы выполнения команд} (в \texttt{E1}), \texttt{Создание/обновление конфигурации и статусов} (в \texttt{D2}), \texttt{Сигналы запуска/паузы/возобновления} (в \texttt{P1}), \texttt{Метрики и логи администрирования} (в \texttt{E3}).
  \end{itemize}
  \item \textbf{P3 Получение/выдача данных.} Обслуживает клиентские запросы на чтение.
  Потоки:
  \begin{itemize}
    \item вход: \texttt{Запросы на получение блокчейн-данных} (из \texttt{E1}), \texttt{Данные для fallback-ответа} (из \texttt{E2});
    \item чтение: \texttt{Чтение индексированных данных} (из \texttt{D1});
    \item выход: \texttt{Ответы с данными/ошибками} (в \texttt{E1}), \texttt{RPC fallback при отсутствии данных} (в \texttt{E2}), \texttt{Дозапись полученных fallback-данных} (в \texttt{D1}), \texttt{Метрики и логи выдачи данных} (в \texttt{E3}).
  \end{itemize}
\end{enumerate}
\noindent На DFD 1 уровня также выделены хранилища: \texttt{D1 Операционное хранилище индексера} и \texttt{D2 Конфигурация и состояние jobs}.

\subsubsection{DFD уровня 2: детализация процесса индексации}
\noindent DFD 2 уровня для процесса \texttt{P1} (приложение В, DFD 2 -- индексация) включает подпроцессы \texttt{P1.1--P1.5}:
\begin{enumerate}
  \item \textbf{P1.1 Получение tip и диапазона синхронизации.}
  Потоки:
  \begin{itemize}
    \item вход: \texttt{Активные jobs и текущий progress\_height} (из \texttt{D2});
    \item выход: \texttt{RPC: getblockcount/getblockhash} (в \texttt{E2}), \texttt{Диапазон высот для обработки} (в \texttt{P1.2});
    \item вход-ответ: \texttt{Tip и контрольные хэши} (из \texttt{E2}).
  \end{itemize}
  \item \textbf{P1.2 Загрузка блока и транзакций.}
  Потоки:
  \begin{itemize}
    \item вход: \texttt{Диапазон высот для обработки} (из \texttt{P1.1});
    \item выход: \texttt{RPC: getblock/getrawtransaction} (в \texttt{E2}), \texttt{Кандидатные блоки и ссылки prev\_hash} (в \texttt{P1.3});
    \item вход-ответ: \texttt{Сырые блоки и транзакции} (из \texttt{E2}).
  \end{itemize}
  \item \textbf{P1.3 Проверка reorg и согласование ветки.}
  Потоки:
  \begin{itemize}
    \item вход: \texttt{Кандидатные блоки и ссылки prev\_hash} (из \texttt{P1.2});
    \item чтение: \texttt{Чтение локальной канонической цепи} (из \texttt{D1});
    \item вход-ответ: \texttt{Локальные hash/height для сверки} (из \texttt{D1});
    \item выход: \texttt{Согласованные данные или команды rollback} (в \texttt{P1.4}).
  \end{itemize}
  \item \textbf{P1.4 Нормализация и запись данных.}
  Потоки:
  \begin{itemize}
    \item вход: \texttt{Согласованные данные или команды rollback} (из \texttt{P1.3});
    \item выход: \texttt{Upsert blocks/transactions/inputs/outputs/utxo} (в \texttt{D1}), \texttt{Факт записи и новый checkpoint} (в \texttt{P1.5}).
  \end{itemize}
  \item \textbf{P1.5 Обновление прогресса job.}
  Потоки:
  \begin{itemize}
    \item вход: \texttt{Факт записи и новый checkpoint} (из \texttt{P1.4});
    \item выход: \texttt{Обновление статуса/progress/last\_error} (в \texttt{D2}).
  \end{itemize}
\end{enumerate}

\subsubsection{DFD уровня 2: детализация процесса администрирования}
\noindent DFD 2 уровня для процесса \texttt{P2} (приложение В, DFD 2 -- администрирование) включает подпроцессы \texttt{P2.1--P2.4}:
\begin{enumerate}
  \item \textbf{P2.1 Прием и аутентификация команды.}
  Потоки:
  \begin{itemize}
    \item вход: \texttt{Команды: create/start/pause/resume/stop/status} (из \texttt{E1});
    \item выход: \texttt{Аутентифицированная команда} (в \texttt{P2.2}).
  \end{itemize}
  \item \textbf{P2.2 Валидация перехода состояния.}
  Потоки:
  \begin{itemize}
    \item вход: \texttt{Аутентифицированная команда} (из \texttt{P2.1});
    \item чтение: \texttt{Чтение текущего состояния job} (из \texttt{D2});
    \item вход-ответ: \texttt{Текущее состояние и политика переходов} (из \texttt{D2});
    \item выход: \texttt{Разрешенная команда или ошибка валидации} (в \texttt{P2.3}).
  \end{itemize}
  \item \textbf{P2.3 Применение команды к job.}
  Потоки:
  \begin{itemize}
    \item вход: \texttt{Разрешенная команда или ошибка валидации} (из \texttt{P2.2});
    \item запись: \texttt{Обновление статуса/конфига/временных меток} (в \texttt{D2});
    \item вход-ответ: \texttt{Подтверждение записи} (из \texttt{D2});
    \item выход: \texttt{Итог выполнения команды} (в \texttt{P2.4}).
  \end{itemize}
  \item \textbf{P2.4 Формирование ответа API.}
  Потоки:
  \begin{itemize}
    \item вход: \texttt{Итог выполнения команды} (из \texttt{P2.3});
    \item выход: \texttt{API-ответ: новый статус/ошибка} (в \texttt{E1}).
  \end{itemize}
\end{enumerate}

\subsubsection{DFD уровня 2: детализация процесса получения данных}
\noindent DFD 2 уровня для процесса \texttt{P3} (приложение В, DFD 2 -- получение данных) включает подпроцессы \texttt{P3.1--P3.5}:
\begin{enumerate}
  \item \textbf{P3.1 Прием запроса и валидация параметров.}
  Потоки:
  \begin{itemize}
    \item вход: \texttt{Запросы данных: block/tx/address/mempool} (из \texttt{E1});
    \item выход: \texttt{Нормализованный поисковый запрос} (в \texttt{P3.2}).
  \end{itemize}
  \item \textbf{P3.2 Поиск данных в индексе.}
  Потоки:
  \begin{itemize}
    \item вход: \texttt{Нормализованный поисковый запрос} (из \texttt{P3.1});
    \item чтение: \texttt{Чтение данных по ключам запроса} (из \texttt{D1});
    \item вход-ответ: \texttt{Данные найдены/данные отсутствуют} (из \texttt{D1});
    \item выход: \texttt{Локальные данные (если найдено)} (в \texttt{P3.5}), \texttt{Сигнал fallback (если не найдено)} (в \texttt{P3.3}).
  \end{itemize}
  \item \textbf{P3.3 Получение данных из узла (fallback).}
  Потоки:
  \begin{itemize}
    \item вход: \texttt{Сигнал fallback (если не найдено)} (из \texttt{P3.2});
    \item выход: \texttt{RPC-запрос при отсутствии данных в индексе} (в \texttt{E2}), \texttt{Сырые данные для нормализации} (в \texttt{P3.4});
    \item вход-ответ: \texttt{Данные узла} (из \texttt{E2}).
  \end{itemize}
  \item \textbf{P3.4 Нормализация и дозапись в индекс.}
  Потоки:
  \begin{itemize}
    \item вход: \texttt{Сырые данные для нормализации} (из \texttt{P3.3});
    \item запись: \texttt{Дозапись/обновление индекса} (в \texttt{D1});
    \item выход: \texttt{Данные после нормализации} (в \texttt{P3.5}).
  \end{itemize}
  \item \textbf{P3.5 Формирование и выдача ответа.}
  Потоки:
  \begin{itemize}
    \item вход: \texttt{Локальные данные (если найдено)} (из \texttt{P3.2}) или \texttt{Данные после нормализации} (из \texttt{P3.4});
    \item выход: \texttt{Ответ API} (в \texttt{E1}).
  \end{itemize}
\end{enumerate}

\subsection{Описание условных обозначений}
\noindent В DFD-диаграммах используются обозначения: внешняя сущность (\texttt{E}) -- источник/получатель данных вне границ системы; процесс (\texttt{P}) -- преобразование данных внутри системы; хранилище данных (\texttt{D}) -- устойчивое место хранения; стрелка потока -- направленная передача данных или управляющего сообщения. Нумерация процессов (\texttt{P1}, \texttt{P2}, \texttt{P3}) отражает декомпозицию от уровня 0 к уровням 1 и 2, а именование потоков задается по смыслу передаваемых данных.

\subsection{Характеристика результата}
\noindent В результате сформирован связанный комплект из пяти DFD-диаграмм: 1 диаграмма уровня 0, 1 диаграмма уровня 1 и 3 диаграммы уровня 2. Комплект покрывает полный эксплуатационный контур продукта: индексацию, администрирование и выдачу данных с fallback-сценарием. Такая структура обеспечивает прозрачность движения данных, разделение ответственности между процессами и основу для трассировки требований в архитектуру и тесты.

\subsection{Оценка качества артефакта}
\noindent Качество DFD-артефакта оценивается по следующим критериям.
\begin{enumerate}
  \item \textbf{Полнота потоков.} Проверяется наличие всех обязательных входных, внутренних и выходных потоков данных для целевых сценариев MVP.
  \item \textbf{Балансировка уровней.} Контролируется корректность декомпозиции: потоки на нижнем уровне не должны противоречить потокам верхнего уровня.
  \item \textbf{Семантическая точность.} Оценивается соответствие наименований процессов и потоков фактическим операциям системы.
  \item \textbf{Непротиворечивость с ER и архитектурой.} Проверяется согласованность DFD с инфологической схемой, API-контрактами и модульной структурой продукта.
  \item \textbf{Проверяемость сценариев.} Оценивается возможность построить тестовые случаи напрямую из зафиксированных потоков и переходов.
  \item \textbf{Эксплуатационная применимость.} Проверяется пригодность диаграмм для диагностики инцидентов и анализа узких мест в продуктивном контуре.
\end{enumerate}

\newpage

\section{Модель процессов функционирования продукта и его составных частей}

\subsection{Место, роль и назначение результата}
\noindent В главе используется процессная модель BPMN 2.0, описывающая сквозной бизнес-процесс функционирования индексера и взаимодействие его составных частей. Модель фиксирует последовательность действий от входного запроса до получения результата, определяет точки ветвления, правила обработки исключений и распределение ответственности между слоями API, управления jobs, индексации, RPC-интеграции, хранения и наблюдаемости. Назначение результата -- обеспечить связность требований глав 1--3 с архитектурой и испытаниями.

\subsection{Изложение артефакта}
\noindent BPMN-артефакт выполнен в формате PlantUML и вынесен в приложение Г в виде двух диаграмм:
\begin{enumerate}
  \item \texttt{bpmn-data-serving-process.puml} -- процесс обработки запроса данных от пользователя;
  \item \texttt{bpmn-indexer-process.puml} -- процесс конфигурации индексации администратором и последующего асинхронного выполнения индексирования.
\end{enumerate}
\noindent Такое разбиение соответствует двум ключевым эксплуатационным сценариям системы: пользовательское чтение данных и административное управление жизненным циклом индексации.

\subsubsection{Диаграмма 1. Запрос данных от пользователя}
\noindent Диаграмма \texttt{bpmn-data-serving-process} отражает сквозной путь пользовательского запроса в слоях \texttt{Клиент/Оператор}, \texttt{API Layer}, \texttt{Data Query Layer}, \texttt{RPC Client}.
\begin{enumerate}
  \item \textbf{Инициация запроса.} Пользователь инициирует запрос данных (\texttt{block/tx/address/mempool}); процесс стартует в клиентской дорожке.
  \item \textbf{Входная валидация.} В \texttt{API Layer} выполняются проверка Basic Auth и валидация параметров.
  \item \textbf{Локальный поиск.} В \texttt{Data Query Layer} выполняется поиск в \texttt{D1}. При наличии данных слой подготавливает результат и передает его в \texttt{API Layer}.
  \item \textbf{Формирование ответа из индекса.} \texttt{API Layer} формирует и возвращает ответ клиенту; процесс завершается.
  \item \textbf{Fallback-ветка.} При отсутствии данных в индексе \texttt{RPC Client} запрашивает данные у Bitcoin Core.
  \item \textbf{Нормализация и дозапись.} Полученные данные нормализуются в \texttt{Data Query Layer}, записываются в \texttt{D1} и передаются в \texttt{API Layer}.
  \item \textbf{Финальный ответ после fallback.} \texttt{API Layer} формирует ответ и возвращает его клиенту. При RPC-ошибке \texttt{API Layer} возвращает ошибку доступности источника. Обе ветки завершаются конечным событием.
\end{enumerate}

\subsubsection{Диаграмма 2. Конфигурация индексации от администратора}
\noindent Диаграмма \texttt{bpmn-indexer-process} отражает административный запуск/изменение состояния jobs и асинхронный цикл индексации в дорожках \texttt{Клиент/Оператор}, \texttt{API Layer}, \texttt{JobService}, \texttt{Indexer Pipeline}, \texttt{RPC Client}, \texttt{Monitoring}.
\begin{enumerate}
  \item \textbf{Инициация административной команды.} Администратор отправляет команду \texttt{start/pause/resume/stop/status}.
  \item \textbf{Проверка и валидация.} В \texttt{API Layer} выполняются аутентификация и валидация структуры команды.
  \item \textbf{Контроль перехода состояния.} \texttt{JobService} проверяет допустимость перехода и при успехе обновляет состояние job в \texttt{D2}.
  \item \textbf{Немедленный ответ администратору.} После принятия валидной команды \texttt{API Layer} сразу возвращает подтверждение принятия (без ожидания завершения индексации). При невалидной команде сразу возвращается ошибка.
  \item \textbf{Запуск асинхронного индексирования.} \texttt{Indexer Pipeline} получает сигнал и выполняет цикл: получение tip, загрузка данных по RPC, проверка reorg, откат при необходимости, запись блока/tx/utxo, обновление \texttt{progress\_height}, ожидание при отсутствии отставания.
  \item \textbf{Наблюдаемость в ходе выполнения.} \texttt{Monitoring} публикует метрики старта run, итерационные метрики (lag, throughput, rpc latency, heartbeat) и терминальные метрики run (status/duration/last\_error).
  \item \textbf{Завершение run.} При переходе \texttt{Job active? = Нет} цикл завершается, фиксируются итоговые метрики/логи и завершается процесс индексирования.
\end{enumerate}

\subsection{Описание условных обозначений}
\noindent В BPMN-модели используются стандартные элементы: событие начала/завершения процесса, задачи (activity), эксклюзивные шлюзы (decision gateway), последовательные потоки (sequence flow), а также swimlane/partition для разделения зон ответственности между компонентами. Нотация соответствует спецификации BPMN 2.0 (OMG) \cite{src_bpmn}.

\subsection{Характеристика результата}
\noindent Сформирован комплект из двух взаимодополняющих BPMN-диаграмм: пользовательского процесса получения данных и административного процесса конфигурации/индексации. Такое представление обеспечивает четкое разделение синхронного клиентского контура (ответ в рамках одного запроса) и асинхронного фонового контура индексирования (длительный run с публикацией метрик). В результате достигаются однозначность ответственности компонентов, трассируемость сценариев и полнота покрытия ключевых эксплуатационных режимов.

\subsection{Оценка качества артефакта}
\noindent Качество BPMN-артефакта оценивается по критериям:
\begin{enumerate}
  \item \textbf{Процессная полнота.} Покрытие всех обязательных ветвей поведения MVP: пользовательское чтение данных, fallback в узел, управление jobs, асинхронная индексация и наблюдаемость.
  \item \textbf{Логическая корректность переходов.} Отсутствие противоречивых ветвлений и недостижимых состояний процесса.
  \item \textbf{Трассируемость к DFD и архитектуре.} Соответствие шагов BPMN потокам данных главы 3 и модульной декомпозиции главы 5.
  \item \textbf{Проверяемость.} Возможность напрямую сформировать тестовые сценарии из узлов и переходов диаграммы.
  \item \textbf{Эксплуатационная применимость.} Пригодность модели для анализа инцидентов, узких мест и регламентов сопровождения.
\end{enumerate}

\newpage

\section{Описание архитектуры программного продукта}

\subsection{Место, роль и назначение результата}
\noindent Результатом главы является архитектурное описание программного продукта в терминах ГОСТ Р 57100--2016 (ISO/IEC/IEEE 42010:2011) \cite{src_gost_57100}. Роль результата -- формально зафиксировать объект архитектурного описания, системный контекст, заинтересованные стороны, состав компонентов, интерфейсы и принципы развития продукта. Назначение результата -- обеспечить согласованность требований глав 1--4 с реализуемой архитектурой MVP и задать основу для контролируемой эволюции решения.

\subsection{Изложение артефакта}
\noindent Архитектурный артефакт включает представление по ГОСТ Р 57100, компонентную архитектуру в нотации C4, формализованное описание интерфейсов взаимодействия, модель развертывания и принципы развития продукта. Такой состав позволяет связать concerns заинтересованных сторон с конкретными архитектурными представлениями и проверяемыми техническими решениями.

\subsubsection{Архитектурное представление в терминах ГОСТ Р 57100}
\noindent Объект архитектурного описания: \textbf{АС «Bitcoin Blockchain Indexer»}.

\noindent Контекст (системное окружение):
\begin{enumerate}
  \item \textbf{Клиентские системы} (CLI, админ-интерфейс, внутренние сервисы) -- источник административных команд и пользовательских запросов данных;
  \item \textbf{Bitcoin Core} (через \texttt{nginx-rpc proxy}) -- внешний источник блоков, транзакций и mempool;
  \item \textbf{PostgreSQL} -- внешнее операционное хранилище бизнес-состояния и индекса;
  \item \textbf{Monitoring Stack} (Prometheus + сборщик логов) -- внешний контур наблюдаемости.
\end{enumerate}

\noindent Архитектурные точки зрения:
\begin{enumerate}
  \item \textbf{Функциональная точка зрения.} Определяет процессы управления jobs, индексирования и выдачи данных, а также их границы ответственности.
  \item \textbf{Информационная точка зрения.} Определяет структуру и жизненный цикл данных (\texttt{blocks}, \texttt{transactions}, \texttt{utxo}, \texttt{jobs}, балансы адресов) в связке с ER-моделью главы 2.
  \item \textbf{Технологическая точка зрения.} Фиксирует выбранный стек и интерфейсы интеграции (Rust, PostgreSQL, JSON-RPC, Prometheus, Docker).
  \item \textbf{Эксплуатационная точка зрения.} Описывает наблюдаемость, управляемость, отказоустойчивость и процедуры восстановления после сбоев.
\end{enumerate}
\noindent В таблице 5.1 приведено соответствие заинтересованных сторон, ключевых concerns и архитектурных представлений, в которых эти concerns отражены.

\begin{table}[H]
\centering
\begin{tabular}{p{0.26\textwidth} p{0.32\textwidth} p{0.34\textwidth}}
\toprule
Заинтересованная сторона & Ключевые интересы (concerns) & Представления/модели, где отражено \\
\midrule
Владелец продукта / заказчик & Прозрачность обработки данных, управляемость статусов jobs, достижение целевых метрик качества & BPMN (приложение с BPMN-диаграммами), DFD (приложение с DFD-диаграммами), C4 (приложение с C4-диаграммой), глава 6 (показатели качества) \\
\midrule
Администратор системы & Надежное управление lifecycle jobs, немедленный ответ API, предсказуемые переходы состояний & BPMN (процесс конфигурации индексации), DFD 1/2 (администрирование), C4: \texttt{API Module} + \texttt{Jobs Module} \\
\midrule
Потребитель API & Стабильный доступ к данным block/tx/address/mempool, корректные fallback-сценарии, приемлемая латентность & BPMN (процесс запроса данных), DFD 1/2 (получение/выдача данных), C4: \texttt{API Module} + \texttt{Storage Module} + \texttt{RPC Module} \\
\midrule
ИТ-эксплуатация & Наблюдаемость, диагностируемость, безопасный перезапуск, контейнеризованное развертывание & C4 (компоненты \texttt{Metrics}/\texttt{Logging}), модель развертывания 5.2.4, глава 7 (испытания) \\
\midrule
Информационная безопасность & Защищенный внешний периметр, контроль доступа, минимизация поверхности атак & C4 и таблицы интерфейсов 5.2.3 (HTTPS, Basic Auth, mTLS), модель развертывания 5.2.4 \\
\bottomrule
\end{tabular}
\caption{Матрица соответствия заинтересованных сторон, их concerns и архитектурных представлений}
\end{table}

\subsubsection{Компонентная архитектура системы}
\noindent Компонентная архитектура представлена в нотации C4 и вынесена в приложение Д (файл \texttt{report/c4/c4-component-architecture.puml}). Диаграмма фиксирует единый контейнер \texttt{Bitcoin Blockchain Indexer} (модульный монолит) и внутренние компоненты:
\begin{enumerate}
  \item \texttt{API Module} -- внешний REST-контур, аутентификация, валидация, маршрутизация;
  \item \texttt{Config Module} -- загрузка и проверка конфигурации, передача параметров модулям;
  \item \texttt{Jobs Module} -- управление заданиями индексации и переходами состояний;
  \item \texttt{Indexer Module} -- core-логика индексации, обработка reorg и mempool;
  \item \texttt{RPC Module} -- безопасная интеграция с Bitcoin Core по JSON-RPC;
  \item \texttt{Storage Module} -- репозитории и транзакционные операции с БД;
  \item \texttt{Metrics Module} -- сбор и экспорт метрик;
  \item \texttt{Logging Module} -- структурированное журналирование.
\end{enumerate}
\noindent Внешние сущности C4: \texttt{Клиентские системы}, \texttt{Bitcoin Core}, \texttt{PostgreSQL}, \texttt{Monitoring Stack}.

\subsubsection{Взаимодействие компонентов}
\noindent Для устранения неоднозначности интерфейсы взаимодействия фиксируются явно для внешнего и внутреннего контуров.

\textbf{Внешние интерфейсы системы представлены в таблице 5.2.}

\begin{table}[H]
\centering
\begin{tabular}{p{0.19\textwidth} p{0.19\textwidth} p{0.17\textwidth} p{0.34\textwidth}}
\toprule
Источник & Получатель & Интерфейс/протокол & Назначение и передаваемые данные \\
\midrule
Клиентские системы & API Module & REST over HTTPS + Basic Auth & Команды \texttt{/v1/jobs/*}, запросы \texttt{block/tx/address/mempool}, получение синхронных API-ответов \\
\midrule
RPC Module & Bitcoin Core & JSON-RPC over HTTPS (mTLS + Basic Auth) & Вызовы \texttt{getblockcount}, \texttt{getblockhash}, \texttt{getblock}, \texttt{getrawtransaction}, получение сырых блокчейн-данных \\
\midrule
Storage Module & PostgreSQL & SQL/TCP (защищенный сетевой сегмент) & Транзакционная запись/чтение \texttt{blocks}, \texttt{transactions}, \texttt{utxo}, \texttt{jobs}, балансов и служебных состояний \\
\midrule
Metrics/Logging Module & Monitoring Stack & Prometheus scrape + JSON log stream & Экспорт технических и бизнес-метрик, публикация структурированных логов для диагностики и аудита \\
\bottomrule
\end{tabular}
\caption{Внешние интерфейсы между индексером и системным окружением}
\end{table}

\textbf{Внутренние интерфейсы компонентов модульного монолита представлены в таблице 5.3.}

\begin{table}[H]
\centering
\begin{tabular}{p{0.19\textwidth} p{0.19\textwidth} p{0.17\textwidth} p{0.34\textwidth}}
\toprule
Источник & Получатель & Интерфейс & Назначение взаимодействия \\
\midrule
API Module & Jobs Module & In-process service call & Передача команд lifecycle jobs; получение результата валидации перехода состояния \\
\midrule
API Module & Storage Module & In-process repository call & Чтение индексированных данных для пользовательских запросов \\
\midrule
Jobs Module & Storage Module & In-process repository call & Чтение/запись статусов jobs, \texttt{progress\_height}, \texttt{last\_error}, таймстемпов выполнения \\
\midrule
Jobs Module & Indexer Module & In-process command signal & Запуск/пауза/возобновление асинхронного цикла индексации \\
\midrule
Indexer Module & RPC Module & In-process client call & Запрос данных blockchain/mempool для этапов синхронизации и fallback \\
\midrule
Indexer Module & Storage Module & In-process repository call & Upsert нормализованных сущностей, rollback при reorg, фиксация checkpoint \\
\midrule
Config Module & API/Jobs/Indexer/RPC & Startup config contract & Передача параметров безопасности, сети, polling, batching, concurrency и стартового набора jobs \\
\bottomrule
\end{tabular}
\caption{Внутренние интерфейсы компонентов модульного монолита}
\end{table}

\noindent Для сценария \textbf{«запрос данных»} основной маршрут: \texttt{Клиент} \textrightarrow{} \texttt{API Module} \textrightarrow{} \texttt{Storage Module}; fallback-маршрут: \texttt{API/Indexer} \textrightarrow{} \texttt{RPC Module} \textrightarrow{} \texttt{Bitcoin Core} с последующей нормализацией и дозаписью.

\noindent Для сценария \textbf{«конфигурация индексации»}: \texttt{Клиент} \textrightarrow{} \texttt{API Module} \textrightarrow{} \texttt{Jobs Module}, после чего API немедленно возвращает ответ о принятии команды, а \texttt{Indexer Module} продолжает асинхронное выполнение.

\subsubsection{Модель развертывания}
\noindent Модель развертывания MVP соответствует контейнерному сценарию \texttt{docker-compose}, вынесена в приложение Д (файл \texttt{report/c4/c4-deployment-architecture.puml}) и включает:
\begin{enumerate}
  \item контейнер \texttt{indexer} (экземпляр модульного монолита, REST API и фоновые workers);
  \item контейнер \texttt{postgres} (персистентное операционное хранилище);
  \item контейнер \texttt{nginx-rpc proxy} (терминация и проксирование защищенного RPC-канала);
  \item внешний узел \texttt{Bitcoin Core};
  \item внешний контур \texttt{Prometheus/log collector}.
\end{enumerate}
\noindent Ключевое архитектурное свойство модели -- stateless-характер экземпляра \texttt{indexer}: состояние jobs, прогресс индексации и бизнес-данные размещены вне процесса. Это обеспечивает безопасный перезапуск экземпляра, переносимость между средами и масштабирование без потери состояния.

\subsubsection{Принципы развития продукта}
\noindent Развитие продукта выполняется по принципам эволюционной архитектуры без нарушения базовой модели модульного монолита:
\begin{enumerate}
  \item \textbf{Эволюция через модули.} Новая функциональность добавляется отдельными модулями с явными контрактами, без размывания ответственности существующих компонентов.
  \item \textbf{Стабильность внешних интерфейсов.} API и RPC-контракты развиваются с обратной совместимостью (версионирование endpoint, контролируемая деприкация).
  \item \textbf{Согласованная эволюция данных.} Изменения схемы БД проводятся только через версионируемые миграции с возможностью контролируемого отката.
  \item \textbf{Наблюдаемость как критерий готовности.} Любое расширение сопровождается метриками и логами, достаточными для диагностики и регрессионного контроля.
  \item \textbf{Безопасность по умолчанию.} Новые интеграции допускаются только при соблюдении политики аутентификации, шифрования каналов и минимально необходимых прав доступа.
  \item \textbf{Пороговая декомпозиция.} Выделение части модулей в отдельные сервисы допускается только при достижении измеримых порогов нагрузки/связности, подтвержденных метриками.
\end{enumerate}

\subsection{Описание условных обозначений}
\noindent В C4-представлении используются обозначения: \texttt{Person} -- внешний актор; \texttt{System\_Ext} -- внешняя система окружения; \texttt{Container\_Boundary} -- граница приложения; \texttt{Component} -- логический компонент внутри контейнера; \texttt{Rel} -- направленная связь (интерфейс) между элементами.

\noindent В таблицах интерфейсов применяются обозначения:
\begin{enumerate}
  \item \textbf{In-process service call} -- вызов между модулями в рамках одного процесса;
  \item \textbf{In-process repository call} -- обращение к слою хранения внутри приложения;
  \item \textbf{REST over HTTPS} -- внешний HTTP API с TLS-шифрованием;
  \item \textbf{JSON-RPC over HTTPS} -- вызовы к узлу Bitcoin Core через защищенный RPC-канал.
\end{enumerate}

\subsection{Характеристика результата}
\noindent Сформировано архитектурное описание, согласованное с предыдущими главами: потоки DFD главы 3, процессы BPMN главы 4 и ER-модель главы 2 трассируются к компонентам C4-модели, их интерфейсам и модели развертывания. Дополнительно зафиксированы принципы развития продукта, что делает архитектуру не только описательной для текущего MVP, но и управляемой для последующих итераций.

\subsection{Оценка качества артефакта}
\noindent Качество архитектурного артефакта оценивается по критериям:
\begin{enumerate}
  \item \textbf{Соответствие стандарту.} Наличие архитектурных представлений и точек зрения в логике ГОСТ Р 57100.
  \item \textbf{Полнота компонентного состава.} Покрытие всех обязательных подсистем MVP.
  \item \textbf{Ясность ответственности.} Отсутствие неоднозначности в распределении функций между компонентами.
  \item \textbf{Полнота интерфейсного описания.} Наличие формализованных внешних и внутренних интерфейсов с указанием протоколов и назначения потоков данных.
  \item \textbf{Корректность взаимодействий.} Соответствие связей компонентам данных, процессам и требованиям безопасности.
  \item \textbf{Трассируемость.} Возможность связать архитектурные решения с требованиями, тестами и эксплуатационными сценариями.
  \item \textbf{Управляемость развития.} Наличие явных принципов эволюции продукта без архитектурной деградации.
\end{enumerate}

\newpage

\section{Описание способа и реализации вычисления показателей качества ПО}

\subsection{Место, роль и назначение результата}
\noindent Результатом главы является формализованный способ определения показателей качества программного продукта «Bitcoin Blockchain Indexer» на этапе MVP. Роль результата -- связать требования к качеству с измеримыми метриками, источниками данных и процедурами оценки. Назначение результата -- обеспечить воспроизводимую основу приемки и регрессионного контроля качества в соответствии с ISO/IEC 25010 (ГОСТ Р ИСО/МЭК 25010--2015) \cite{src_gost_25010}.

\subsection{Изложение артефакта}
\noindent Артефакт включает: (1) выбор характеристик качества по ISO/IEC 25010, релевантных архитектуре модульного монолита индексера; (2) показатели и способы измерения с целевыми значениями MVP; (3) методику определения показателей и описание источников тестовых данных.

\subsubsection{Показатели качества АС по ISO/IEC 25010 (ГОСТ Р ИСО/МЭК 25010-2015)}
\begin{table}[H]
\centering
\begin{tabular}{p{0.20\textwidth} p{0.20\textwidth} p{0.34\textwidth} p{0.18\textwidth}}
\toprule
Характеристика качества & Подхарактеристика & Показатель / способ измерения & Целевое значение (этап MVP) \\
\midrule
Функциональная пригодность & Функциональная полнота & Доля обязательных сценариев API и jobs-lifecycle, подтвержденных тестами: $K_{cov}=\frac{N_{ok}}{N_{req}}$ & $K_{cov} \ge 0.95$ \\
\midrule
Функциональная пригодность & Функциональная корректность & Доля корректных ответов и консистентных записей БД по инвариантам PK/FK/UQ: $K_{corr}=\frac{N_{corr}}{N_{all}}$ & $K_{corr}=1.0$ для приемочного набора \\
\midrule
Надежность & Доступность & Доля времени готовности API: $A=\frac{T_{up}}{T_{total}}$ по метрикам uptime/healthcheck & $A \ge 0.99$ в окне испытаний \\
\midrule
Надежность & Восстанавливаемость & Время восстановления после перезапуска с сохранением progress jobs и без потери данных & $T_{rec} \le 60$ с \\
\midrule
Производительность & Временное поведение & Перцентиль задержки API чтения (\texttt{p95}) по профилю \texttt{normal} & $L_{api,p95} \le 500$ мс \\
\midrule
Производительность & Пропускная способность & Количество обработанных транзакций в секунду: $Q=\frac{N_{tx}}{T}$ для индексации на фиксированном диапазоне блоков & Не ниже базовой линии стенда \\
\midrule
Безопасность & Конфиденциальность/целостность канала & Доля внешних интерфейсов с шифрованием и аутентификацией (HTTPS, Basic Auth, mTLS) & $K_{sec}=1.0$ \\
\midrule
Сопровождаемость & Тестируемость & Доля успешно пройденных автотестов при контрольном прогоне CI: $K_{test}=\frac{N_{pass}}{N_{run}}$ & $K_{test}=1.0$ \\
\midrule
Совместимость & Интероперабельность & Успешность взаимодействия с Bitcoin Core JSON-RPC и Monitoring Stack без ошибок протокола & $K_{interop} \ge 0.99$ \\
\bottomrule
\end{tabular}
\caption{Показатели качества индексера по ISO/IEC 25010 для этапа MVP}
\end{table}

\subsubsection{Формализация вычисления показателей}
\noindent Для количественной оценки используются следующие базовые расчеты:
\begin{enumerate}
  \item \textbf{Покрытие требований:} $K_{cov}=\frac{N_{ok}}{N_{req}}$, где $N_{req}$ -- число обязательных требований/сценариев MVP, $N_{ok}$ -- число требований, подтвержденных испытаниями.
  \item \textbf{Корректность обработки:} $K_{corr}=\frac{N_{corr}}{N_{all}}$, где $N_{all}$ -- общее число проверок ответов и записей в БД, $N_{corr}$ -- число проверок без отклонений.
  \item \textbf{Доступность:} $A=\frac{T_{up}}{T_{total}}$, где $T_{up}$ -- время доступности API в окне наблюдения $T_{total}$.
  \item \textbf{Пропускная способность индексации:} $Q=\frac{N_{tx}}{T}$, где $N_{tx}$ -- число успешно индексированных транзакций за интервал $T$.
  \item \textbf{Латентность API:} используется эмпирическое распределение задержек с контролем $L_{api,p95}$ и $L_{api,p99}$.
\end{enumerate}
\noindent Указанные расчеты выполняются на фиксированном стенде и идентичной конфигурации (\texttt{rpc}, \texttt{indexer}, \texttt{jobs}) для обеспечения сопоставимости результатов между прогонами.

\subsubsection{Способы определения показателей и источники тестовых данных}
\noindent Источники тестовых данных и сценарии измерения формируются в связке с архитектурой и процессами из глав 3--5.
\begin{enumerate}
  \item \textbf{Данные блокчейна и транзакций.} Используются ответы Bitcoin Core RPC (\texttt{getblockcount}, \texttt{getblockhash}, \texttt{getblock}, \texttt{getrawtransaction}) на фиксированном диапазоне высот.
  \item \textbf{Синтетические данные отказов/граничных условий.} Формируются сценарии reorg в пределах \texttt{reorg\_depth}, задержки RPC и частичные ошибки источника для проверки recoverability и корректности fallback.
  \item \textbf{Нагрузочные выборки API.} Генерируются профили \texttt{smoke}/\texttt{normal}/\texttt{stress} для endpoint чтения данных и команд lifecycle jobs.
  \item \textbf{Эксплуатационные телеметрические данные.} Используются метрики Prometheus и структурированные JSON-логи для расчета латентности, доступности, throughput и диагностики ошибок \cite{src_prometheus}.
\end{enumerate}

\noindent Методика измерений:
\begin{enumerate}
  \item \textbf{Подготовка стенда.} Фиксируются версия приложения, конфигурация Docker/Compose, параметры БД и RPC, набор jobs.
  \item \textbf{Прогрев и калибровка.} Выполняется прогревочный прогон без фиксации результата для стабилизации кэшей и соединений.
  \item \textbf{Основной прогон.} Выполняются функциональные и нагрузочные сценарии с одновременным сбором метрик/логов и контрольных SQL-проверок целостности.
  \item \textbf{Расчет показателей.} По собранным данным вычисляются $K_{cov}$, $K_{corr}$, $A$, $Q$, $L_{api,p95}$, $L_{api,p99}$ и показатели безопасности/интероперабельности.
  \item \textbf{Сравнение с целями MVP.} Значения сопоставляются с таблицей 6.1; отклонения оформляются как дефекты или ограничения текущего релиза.
\end{enumerate}

\subsection{Описание условных обозначений}
\noindent Используются обозначения: $N_{req}$ -- число обязательных требований; $N_{ok}$ -- число подтвержденных требований; $N_{all}$ -- общее число проверок; $N_{corr}$ -- число корректных проверок; $T_{up}$ -- время доступности API; $T_{total}$ -- полное окно наблюдения; $N_{tx}$ -- число индексированных транзакций; $Q$ -- пропускная способность; $L_{api,p95}$, $L_{api,p99}$ -- перцентили латентности API; $K_{sec}$ -- доля защищенных интерфейсов; $K_{interop}$ -- доля успешных интеграционных вызовов.

\subsection{Характеристика результата}
\noindent В результате сформирована целостная модель оценки качества ПО: выбран набор показателей по ISO/IEC 25010, заданы формулы и правила измерения, определены источники тестовых данных и целевые значения MVP. Модель обеспечивает сопоставимость результатов между прогонами и напрямую поддерживает приемочные решения по готовности версии.

\subsection{Оценка качества артефакта}
\noindent Качество артефакта оценки показателей определяется по критериям:
\begin{enumerate}
  \item \textbf{Стандартная обоснованность.} Соответствие выбранных характеристик и подхарактеристик ISO/IEC 25010.
  \item \textbf{Измеримость.} Наличие формализованных метрик, формул и источников данных для каждого показателя.
  \item \textbf{Воспроизводимость.} Возможность повторить измерения на идентичной конфигурации и получить сопоставимые результаты.
  \item \textbf{Трассируемость.} Связь показателей качества с архитектурой, процессами и испытаниями (главы 3--7).
  \item \textbf{Практическая применимость.} Пригодность метрик для принятия решений о выпуске MVP и для регрессионного контроля.
\end{enumerate}

\newpage

\section{Программа и методики испытаний ПО}

\subsection{Место, роль и назначение результата}
\noindent Результатом главы является программа и методики испытаний прототипа АС ончейн-индексации Bitcoin. Роль результата -- определить регламент проверки функциональных, эксплуатационных и качественных требований MVP, сформулированных в главах 1--6. Назначение результата -- обеспечить воспроизводимую процедуру подтверждения готовности ПО к приемке и последующему развитию.

\subsection{Изложение артефакта}
\subsubsection{Объект испытания}
\noindent Объектом испытаний является прототип АС ончейн-индексации \texttt{Bitcoin Blockchain Indexer}, включающий:
\begin{enumerate}
  \item модульный монолитный сервис (\texttt{API Module}, \texttt{Jobs Module}, \texttt{Indexer Module}, \texttt{RPC Module}, \texttt{Storage Module}, \texttt{Metrics/Logging});
  \item внешний RPC-контур (\texttt{nginx-rpc proxy} + узел Bitcoin Core);
  \item операционное хранилище PostgreSQL;
  \item контур наблюдаемости (Prometheus и сбор структурированных логов);
  \item конфигурацию jobs и механизм lifecycle-команд (\texttt{start/pause/resume/stop/status}).
\end{enumerate}

\subsubsection{Цели испытания}
\noindent Испытания проводятся для достижения следующих целей:
\begin{enumerate}
  \item подтвердить функциональную корректность API и жизненного цикла jobs;
  \item подтвердить корректность индексации блоков/транзакций/UTXO с учетом reorg и fallback-сценариев;
  \item оценить производительность и устойчивость в профилях \texttt{smoke}, \texttt{normal}, \texttt{stress};
  \item подтвердить выполнение целевых показателей качества MVP, определенных в главе 6;
  \item проверить пригодность артефактов наблюдаемости для диагностики отказов и регрессий.
\end{enumerate}

\subsubsection{Виды испытаний}
\begin{table}[H]
\centering
\begin{tabular}{p{0.22\textwidth} p{0.36\textwidth} p{0.32\textwidth}}
\toprule
Вид испытаний & Объект проверки & Критерий успешности \\
\midrule
Модульные тесты & Бизнес-правила модулей конфигурации, jobs, RPC и хранения & Все тесты проходят; критические инварианты не нарушены \\
\midrule
Интеграционные тесты & Связка API--Jobs--Storage--RPC на тестовом контуре & Корректные статусы/ответы API; отсутствие ошибок целостности БД \\
\midrule
Функциональные API-испытания & Endpoint управления jobs и endpoint чтения данных & Коды ответов и тело ответа соответствуют спецификации сценария \\
\midrule
Нагрузочные испытания & Латентность и пропускная способность при конкурентных запросах & Выполнение целевых порогов из главы 6 (\texttt{p95}, throughput) \\
\midrule
Испытания отказоустойчивости & Поведение при сбоях RPC/БД, перезапуске сервиса, временной недоступности узла & Восстановление в пределах целевого времени; отсутствие потери состояния jobs \\
\midrule
Испытания безопасности канала & HTTPS, Basic Auth, mTLS в контуре внешних интерфейсов & Все внешние вызовы выполняются только по защищенным каналам \\
\bottomrule
\end{tabular}
\caption{Виды испытаний, объект проверки и критерии успешности}
\end{table}

\subsubsection{Методики испытаний}
\noindent Для каждого вида испытаний применяется следующая методика.
\begin{enumerate}
  \item \textbf{Модульное тестирование.} Инструмент: \texttt{cargo test}. Порядок: запуск набора unit-тестов по модулям, анализ падений, локализация дефекта до функции/контракта. Проверяется корректность переходов состояний, валидация конфигурации, сериализация RPC, базовые репозитории хранения.
  \item \textbf{Интеграционное тестирование.} Инструменты: \texttt{cargo test} + \texttt{testcontainers} \cite{src_testcontainers}. Порядок: подъем изолированного контейнера БД, запуск сервисных сценариев API и jobs, сверка результата с состоянием БД. Проверяется целостность записей и согласованность межмодульных контрактов.
  \item \textbf{Функциональное API-тестирование.} Инструменты: HTTP-клиент (\texttt{curl}/коллекция запросов), SQL-проверки. Порядок: выполнение позитивных и негативных сценариев endpoint, верификация кодов ответов, проверка бизнес-эффекта в БД и логах.
  \item \textbf{Нагрузочное тестирование.} Инструменты: генератор конкурентных HTTP-запросов, метрики Prometheus. Порядок: профили \texttt{smoke} \textrightarrow{} \texttt{normal} \textrightarrow{} \texttt{stress}, сбор latency/throughput/error-rate. Проверяется устойчивость SLA-метрик на фиксированном тестовом наборе.
  \item \textbf{Испытания отказоустойчивости.} Инструменты: управляемый перезапуск контейнеров Docker/Compose \cite{src_docker,src_compose}, анализ метрик и логов. Порядок: моделирование отказа RPC/БД/сети, возврат к штатному режиму, проверка восстановления прогресса индексации и отсутствия дубликатов.
\end{enumerate}

\subsubsection{Тестовые данные и порядок их подготовки}
\noindent Основным источником тестовых данных является \textbf{засинхронизированная нода Bitcoin Testnet}. Использование Testnet выбрано по следующим причинам:
\begin{enumerate}
  \item существенно меньший объем данных по сравнению с Mainnet снижает требования к диску и времени первичной синхронизации;
  \item меньшие затраты на вычислительные ресурсы и сетевой трафик для регулярных прогонов;
  \item сохранение реалистичной структуры данных блокчейна, достаточной для проверки логики индексации, reorg и mempool.
\end{enumerate}
\noindent Порядок подготовки данных:
\begin{enumerate}
  \item синхронизация Testnet-ноды до фиксированной высоты $H_{base}$;
  \item фиксация диапазона испытаний $[H_{start}, H_{end}]$ и набора контрольных tx/address;
  \item очистка/инициализация тестовой БД и применение миграций;
  \item запуск jobs в режимах \texttt{full-sync} и \texttt{address-list};
  \item сохранение baseline-результатов для последующих регрессионных сравнений.
\end{enumerate}

\subsubsection{Требования к тестовой среде}
\noindent Тестовая среда должна соответствовать следующим требованиям:
\begin{enumerate}
  \item контейнерный запуск через Docker/Compose с выделенными сервисами \texttt{indexer}, \texttt{postgres}, \texttt{nginx-rpc-proxy};
  \item доступность засинхронизированной Testnet-ноды Bitcoin Core;
  \item изолированный экземпляр PostgreSQL для испытаний;
  \item включенные сбор метрик Prometheus и структурированное логирование;
  \item фиксированная конфигурация таймаутов RPC, reorg depth, polling и ограничений параллелизма;
  \item синхронизированное системное время на узлах стенда для корректного расчета latency/availability.
\end{enumerate}

\subsubsection{Порядок проведения испытаний}
\noindent Проведение испытаний выполняется по этапам:
\begin{enumerate}
  \item \textbf{Подготовительный этап:} проверка конфигурации, доступности Testnet-ноды, запуск тестового контура и миграций.
  \item \textbf{Базовый этап:} модульные и интеграционные тесты, подтверждение корректного старта и базовых API-сценариев.
  \item \textbf{Функциональный этап:} проверка полного цикла jobs и сценариев чтения данных, включая fallback при отсутствии данных в индексе.
  \item \textbf{Нагрузочный этап:} последовательный прогон профилей \texttt{smoke/normal/stress} с фиксацией метрик.
  \item \textbf{Отказоустойчивость и восстановление:} моделирование отказов и проверка восстановления состояния/прогресса.
  \item \textbf{Финальный этап:} расчет показателей из главы 6, оформление протокола испытаний и итогового статуса (\texttt{pass/fail}) по каждому виду испытаний.
\end{enumerate}

\subsection{Описание условных обозначений}
\noindent В главе используются обозначения: \texttt{pass} -- испытание пройдено; \texttt{fail} -- испытание не пройдено; \texttt{N/A} -- испытание неприменимо; $H_{start}, H_{end}$ -- границы тестового диапазона высот; $H_{base}$ -- базовая синхронизированная высота Testnet-ноды; \texttt{p95/p99} -- перцентили латентности API.

\subsection{Характеристика результата}
\noindent Сформирована полная программа и методики испытаний прототипа АС ончейн-индексации, включающая объект испытаний, цели, виды испытаний с критериями успешности, методики выполнения, требования к тестовой среде и пошаговый регламент проведения. Результат согласован с архитектурой (глава 5), моделью качества (глава 6) и ранее описанными процессами DFD/BPMN.

\subsection{Оценка качества артефакта}
\noindent Качество программы и методик испытаний оценивается по критериям:
\begin{enumerate}
  \item \textbf{Полнота покрытия.} Покрытие всех критичных сценариев MVP: lifecycle jobs, индексация, fallback, reorg, наблюдаемость.
  \item \textbf{Измеримость.} Наличие однозначных критериев успешности и привязки к метрикам качества главы 6.
  \item \textbf{Воспроизводимость.} Возможность повторного прогона на идентичном стенде с сопоставимыми результатами.
  \item \textbf{Трассируемость.} Связь методик испытаний с требованиями, архитектурой и интерфейсами системы.
  \item \textbf{Эксплуатационная применимость.} Пригодность методик для регулярного регрессионного контроля и приемки релизов.
\end{enumerate}

\newpage

\section*{Заключение}
\addcontentsline{toc}{section}{Заключение}
\noindent В ходе выполнения работы была проработана предметная область ончейн-индексации блокчейна Bitcoin и сформировано проектное решение прототипа автоматизированной системы \texttt{Bitcoin Blockchain Indexer}, объединяющей контур получения данных из узла Bitcoin Core, контур нормализации и хранения данных в PostgreSQL, контур административного управления jobs и контур предоставления данных через REST API.

\noindent По результатам работы получены следующие основные результаты:
\begin{enumerate}
  \item выполнено вербальное описание предметной области и определены роли участников процесса (клиентские системы, администратор, контур RPC-узла, контур мониторинга);
  \item построена инфологическая модель данных (ER-диаграмма в нотации Чена), отражающая сущности блоков, транзакций, входов/выходов, UTXO, заданий индексации и балансов;
  \item разработана модель функционирования системы в виде DFD-диаграмм уровней 0--2, определяющая границы АС, процессы индексации/администрирования/выдачи данных и хранилища;
  \item описаны процессы функционирования продукта в BPMN для двух ключевых сценариев: запрос данных пользователем и конфигурация индексации администратором с асинхронным выполнением;
  \item сформировано архитектурное описание по ГОСТ Р 57100--2016, включая C4-компонентную модель модульного монолита, интерфейсы взаимодействия и модель развертывания;
  \item определены показатели качества ПО по ISO/IEC 25010 (ГОСТ Р ИСО/МЭК 25010--2015), способы их измерения и целевые значения для этапа MVP;
  \item разработаны программа и методики испытаний, включающие объект и цели испытаний, виды и порядок испытаний, требования к тестовой среде и подготовку тестовых данных на базе засинхронизированной Testnet-ноды.
\end{enumerate}

\noindent Практическая часть подтверждена проверками компилируемости, модульного и интеграционного тестирования, а также контейнерной сборкой. Итогом работы является согласованный комплект проектных и испытательных артефактов, подтверждающий достижение целевого уровня качества MVP и готовность к дальнейшему развитию функциональности (углубленная индексация, расширенные режимы обработки reorg/mempool, развитие эксплуатационной аналитики и масштабирования).

\newpage

\section*{СПИСОК ИСПОЛЬЗОВАННЫХ ИСТОЧНИКОВ}
\addcontentsline{toc}{section}{СПИСОК ИСПОЛЬЗОВАННЫХ ИСТОЧНИКОВ}
\begin{thebibliography}{99}
\bibitem{src_rust}
Клабник С., Николс К. Язык программирования Rust [Электронный ресурс]. -- URL: \url{https://doc.rust-lang.org/book/} (дата обращения: 21.02.2026).

\bibitem{src_postgres}
PostgreSQL 16 Documentation [Электронный ресурс]. -- URL: \url{https://www.postgresql.org/docs/16/index.html} (дата обращения: 21.02.2026).

\bibitem{src_bitcoin_core}
Bitcoin Core Developer Documentation [Электронный ресурс]. -- URL: \url{https://developer.bitcoin.org/reference/rpc/index.html} (дата обращения: 21.02.2026).

\bibitem{src_chen}
Chen P. The Entity-Relationship Model -- Toward a Unified View of Data // ACM Transactions on Database Systems. -- 1976. -- Vol. 1, No. 1. -- P. 9--36.

\bibitem{src_bpmn}
Business Process Model and Notation (BPMN) Version 2.0.2 [Электронный ресурс]. -- URL: \url{https://www.omg.org/spec/BPMN/2.0.2} (дата обращения: 21.02.2026).

\bibitem{src_uml}
Unified Modeling Language (UML), Version 2.5.1 [Электронный ресурс]. -- URL: \url{https://www.omg.org/spec/UML/2.5.1} (дата обращения: 21.02.2026).

\bibitem{src_prometheus}
Prometheus Documentation [Электронный ресурс]. -- URL: \url{https://prometheus.io/docs/introduction/overview/} (дата обращения: 21.02.2026).

\bibitem{src_testcontainers}
Testcontainers for Rust Documentation [Электронный ресурс]. -- URL: \url{https://docs.rs/testcontainers/} (дата обращения: 21.02.2026).

\bibitem{src_docker}
Docker Documentation [Электронный ресурс]. -- URL: \url{https://docs.docker.com/} (дата обращения: 21.02.2026).

\bibitem{src_compose}
Docker Compose Documentation [Электронный ресурс]. -- URL: \url{https://docs.docker.com/compose/} (дата обращения: 21.02.2026).

\bibitem{src_gost}
ГОСТ 7.32--2017. Система стандартов по информации, библиотечному и издательскому делу. Отчет о научно-исследовательской работе. Структура и правила оформления. -- М.: Стандартинформ, 2017. -- 27 с.

\bibitem{src_gost_57100}
ГОСТ Р 57100--2016 (ISO/IEC/IEEE 42010:2011). Системная и программная инженерия. Описание архитектуры. -- М.: Стандартинформ, 2016.

\bibitem{src_gost_25010}
ГОСТ Р ИСО/МЭК 25010--2015. Системная и программная инженерия. Требования и оценка качества систем и программного обеспечения (SQuaRE). Модели качества систем и программных продуктов. -- М.: Стандартинформ, 2015.
\end{thebibliography}

\newpage

\appendix
\section{Пример конфигурации YAML}
\begin{lstlisting}
server:
  bind_host: "0.0.0.0"
  bind_port: 8443
  tls:
    cert_path: "/etc/indexer/tls/server.crt"
    key_path: "/etc/indexer/tls/server.key"
  auth:
    basic:
      username: "admin"
      password_env: "INDEXER_API_PASSWORD"

rpc:
  node_id: "btc-mainnet-1"
  url: "https://nginx-rpc:443"
  auth:
    basic:
      username: "rpcuser"
      password_env: "BITCOIN_RPC_PASSWORD"
  mtls:
    enabled: true
    ca_path: "/etc/indexer/mtls/ca.crt"
    client_cert_path: "/etc/indexer/mtls/client.crt"
    client_key_path: "/etc/indexer/mtls/client.key"
  timeouts:
    connect_ms: 5000
    request_ms: 30000

indexer:
  chain: "bitcoin"
  network: "mainnet"
  reorg_depth: 12
  poll:
    tip_interval_ms: 5000
    mempool_interval_ms: 3000
  concurrency:
    max_jobs: 5
    rpc_parallelism: 8
    db_writer_parallelism: 4
  batching:
    blocks_per_batch: 50
    txs_per_batch: 5000

jobs:
  - job_id: "full-sync"
    mode: "all_addresses"
    enabled: true

  - job_id: "watchlist"
    mode: "address_list"
    enabled: false
    addresses:
      - "bc1q...."
      - "1A1zP1eP5QGefi2DMPTfTL5SLmv7DivfNa"
\end{lstlisting}

\newpage
\section{ER-диаграмма предметной области (нотация Чена)}
\begin{figure}[H]
\centering
\resizebox{\textwidth}{!}{%
\begin{tikzpicture}[node distance=1.0cm and 1.2cm, auto, >=Stealth, every node/.style={font=\scriptsize}]
  \tikzstyle{ent}=[draw, rectangle, minimum width=2.8cm, minimum height=0.9cm, align=center]
  \tikzstyle{rel}=[draw, diamond, aspect=2.0, minimum width=2.3cm, minimum height=0.8cm, align=center]
  \tikzstyle{attr}=[draw, ellipse, minimum width=2.3cm, minimum height=0.7cm, align=center]

  \node[ent] (blocks) {blocks};
  \node[ent, right=3.2cm of blocks] (txs) {transactions};
  \node[ent, right=3.0cm of txs] (outs) {tx\_outputs};
  \node[ent, below=2.1cm of txs] (ins) {tx\_inputs};
  \node[ent, below=2.1cm of outs] (utxo) {utxos\_current};
  \node[ent, below=2.1cm of blocks] (jobs) {jobs};
  \node[ent, below=2.1cm of jobs] (jobaddr) {job\_addresses};
  \node[ent, right=3.0cm of utxo] (balcur) {address\_balance\_current};
  \node[ent, below=2.1cm of balcur] (balhist) {address\_balance\_history};
  \node[ent, right=3.0cm of jobs] (health) {node\_health};

  \node[rel] (r1) at ($(blocks)!0.5!(txs)$) {contains};
  \node[rel] (r2) at ($(txs)!0.5!(outs)$) {has\_output};
  \node[rel] (r3) at ($(txs)!0.5!(ins)$) {has\_input};
  \node[rel] (r4) at ($(ins)!0.5!(outs)$) {spends};
  \node[rel] (r5) at ($(outs)!0.5!(utxo)$) {forms\_utxo};
  \node[rel] (r6) at ($(utxo)!0.5!(balcur)$) {updates\_balance};
  \node[rel] (r7) at ($(balcur)!0.5!(balhist)$) {snapshots};
  \node[rel] (r8) at ($(jobs)!0.5!(jobaddr)$) {has\_filter};
  \node[rel] (r9) at ($(jobs)!0.5!(health)$) {monitors\_node};

  \draw (blocks) -- node[above]{1} (r1);
  \draw (r1) -- node[above]{N} (txs);
  \draw (txs) -- node[above]{1} (r2);
  \draw (r2) -- node[above]{N} (outs);
  \draw (txs) -- node[right]{1} (r3);
  \draw (r3) -- node[left]{N} (ins);
  \draw (ins) -- node[above]{N} (r4);
  \draw (r4) -- node[above]{1} (outs);
  \draw (outs) -- node[right]{1} (r5);
  \draw (r5) -- node[left]{0..1} (utxo);
  \draw (utxo) -- node[above]{N} (r6);
  \draw (r6) -- node[above]{1} (balcur);
  \draw (balcur) -- node[right]{1} (r7);
  \draw (r7) -- node[left]{N} (balhist);
  \draw (jobs) -- node[left]{1} (r8);
  \draw (r8) -- node[right]{N} (jobaddr);
  \draw (jobs) -- node[above]{N} (r9);
  \draw (r9) -- node[above]{1} (health);

  \node[attr, above=0.9cm of blocks] (a1) {\underline{hash}};
  \node[attr, above=0.9cm of txs] (a2) {\underline{txid}};
  \node[attr, above=0.9cm of outs] (a3) {\underline{(txid, vout)}};
  \node[attr, below=0.9cm of ins] (a4) {\underline{(txid, vin)}};
  \node[attr, right=0.9cm of utxo] (a5) {\underline{(out\_txid, out\_vout)}};
  \node[attr, left=0.9cm of jobs] (a6) {\underline{job\_id}};
  \node[attr, left=0.9cm of jobaddr] (a7) {\underline{(job\_id, address)}};
  \node[attr, right=0.9cm of balcur] (a8) {\underline{address}};
  \node[attr, right=0.9cm of balhist] (a9) {\underline{(address, block\_height)}};
  \node[attr, right=0.9cm of health] (a10) {\underline{node\_id}};

  \draw (blocks) -- (a1);
  \draw (txs) -- (a2);
  \draw (outs) -- (a3);
  \draw (ins) -- (a4);
  \draw (utxo) -- (a5);
  \draw (jobs) -- (a6);
  \draw (jobaddr) -- (a7);
  \draw (balcur) -- (a8);
  \draw (balhist) -- (a9);
  \draw (health) -- (a10);
\end{tikzpicture}
}
\caption{ER-диаграмма предметной области индексера в нотации Чена}
\label{fig:chen-er}
\end{figure}

\newpage
\section{Набор DFD-диаграмм для модели внедрения}
\noindent В данном приложении используются отдельные исходные файлы PlantUML, по которым формируются изображения DFD-диаграмм:
\begin{enumerate}
  \item \texttt{report/dfd/dfd-level-0-context.puml} -- DFD уровня 0 (контекстная диаграмма);
  \item \texttt{report/dfd/dfd-level-1-decomposition.puml} -- DFD уровня 1;
  \item \texttt{report/dfd/dfd-level-2-indexing.puml} -- DFD уровня 2, детализация процесса индексации;
  \item \texttt{report/dfd/dfd-level-2-admin.puml} -- DFD уровня 2, детализация процесса администрирования;
  \item \texttt{report/dfd/dfd-level-2-data-serving.puml} -- DFD уровня 2, детализация процесса получения данных.
\end{enumerate}
\noindent Подготовленные изображения диаграмм включаются в итоговую версию пояснительной записки в соответствии с требованиями оформления приложений.

\newpage
\section{BPMN-диаграмма процессов функционирования продукта}
\noindent Для главы 4 используется набор исходных файлов BPMN-диаграмм в формате PlantUML:
\begin{enumerate}
  \item \texttt{report/bpmn/bpmn-data-serving-process.puml} -- BPMN-диаграмма процесса запроса данных от пользователя;
  \item \texttt{report/bpmn/bpmn-indexer-process.puml} -- BPMN-диаграмма процесса конфигурации индексации от администратора и асинхронного выполнения индексации.
\end{enumerate}

\noindent Инструкция по компиляции диаграмм:
\begin{verbatim}
cd report/bpmn
plantuml bpmn-data-serving-process.puml
plantuml bpmn-indexer-process.puml
\end{verbatim}
\noindent Для альтернативных форматов:
\begin{verbatim}
plantuml -tsvg bpmn-data-serving-process.puml
plantuml -tsvg bpmn-indexer-process.puml
plantuml -tpdf bpmn-data-serving-process.puml
plantuml -tpdf bpmn-indexer-process.puml
\end{verbatim}

\newpage
\section{C4-диаграммы архитектуры}
\noindent Для главы 5 используются исходные файлы C4-диаграмм в формате PlantUML:
\begin{enumerate}
  \item \texttt{report/c4/c4-component-architecture.puml} -- компонентная архитектура модульного монолита индексера и внешних систем.
  \item \texttt{report/c4/c4-deployment-architecture.puml} -- модель развертывания MVP в контейнерном окружении.
\end{enumerate}

\noindent Инструкция по компиляции диаграмм:
\begin{verbatim}
cd report/c4
plantuml c4-component-architecture.puml
plantuml c4-deployment-architecture.puml
\end{verbatim}
\noindent Для альтернативных форматов:
\begin{verbatim}
plantuml -tsvg c4-component-architecture.puml
plantuml -tsvg c4-deployment-architecture.puml
plantuml -tpdf c4-component-architecture.puml
plantuml -tpdf c4-deployment-architecture.puml
\end{verbatim}
\noindent Примечание: диаграмма использует \texttt{!includeurl} для C4-PlantUML, поэтому при компиляции требуется доступ в интернет, либо локальная замена на \texttt{!include}.

\end{document}
